\documentclass[11pt]{article}
\usepackage[margin=1.1in]{geometry}
\usepackage{amsmath,amssymb,amsthm,booktabs,hyperref}

\newtheorem{theorem}{Theorem}
\newtheorem{proposition}[theorem]{Proposition}
\newtheorem{lemma}[theorem]{Lemma}
\newtheorem{corollary}[theorem]{Corollary}
\theoremstyle{definition}
\newtheorem{definition}[theorem]{Definition}
\newtheorem{remark}[theorem]{Remark}

\newcommand{\Sn}{\mathfrak{S}_n}
\newcommand{\ars}{a_{r,s}}
\newcommand{\brs}{b_{r,s}}

\title{Holonomicity of the restricted permutation counts $a_{r,s}$ and $b_{r,s}$}
\author{Jaideep Sai Padhi\\ \small jpadhi@purdue.edu}
\date{August 2026}

\begin{document}
\maketitle

\begin{abstract}
For integers $r,s\ge 1$ let $\ars(n)$ count permutations $\pi\in\Sn$ with
$\pi_{i+r}-\pi_i\neq s$ for all $i$, and let $\brs(n)$ count those with
$|\pi_{i+r}-\pi_i|\neq s$. Spahn and Zeilberger \cite{SZ} ask whether these
sequences are holonomic (P-recursive) for all $r,s>1$, noting that no general
theory was known. We answer this affirmatively. The proof generalises a
bijection of Matsuo, used by Spahn and Zeilberger only in the case
$(r,s)=(2,2)$: interleaving the $r$ arithmetic progressions of step $r$ on
positions and the $s$ progressions of step $s$ on values converts the
$(r,s)$-condition into the adjacency condition $\sigma_{k+1}-\sigma_k\neq 1$,
excused on a set of \emph{bounded} size --- exactly $(r-1)+(s-1)$ elements for
$\ars$ and $(r-1)+2(s-1)$ for $\brs$, independently of $n$. Boundedness is the
whole point: it makes the associated inclusion--exclusion representation carry
finitely many catalytic variables for each fixed $(r,s)$, whence holonomicity
follows from standard closure theorems.
\end{abstract}

\section{Introduction}

Let $r,s\ge 1$ and put
\[
  \ars(n) = \#\{\pi\in\Sn : \pi_{i+r}-\pi_i \neq s \ \ (1\le i\le n-r)\},
  \qquad
  \brs(n) = \#\{\pi\in\Sn : |\pi_{i+r}-\pi_i| \neq s\}.
\]
For $(r,s)=(1,1)$ these are classical: $a_{1,1}$ is counted by Riordan's
recurrence and $b_{1,1}$ is the Hertzsprung problem, OEIS \texttt{A002464}.
For larger $r,s$ the situation is different. Kauers and Koutschan
\cite{KK} guessed an order-$8$, degree-$11$ recurrence for
$a_{2,2}=\texttt{A189281}$ from $35$ terms and remark that they have not tried
to prove it; Spahn and Zeilberger \cite{SZ} pose as their third challenge the
question of whether \emph{all} $\ars,\brs$ with $r,s>1$ are holonomic, observing
that there is no a priori guarantee.

The obstruction is not the existence of a combinatorial decomposition --- one
is available for every $(r,s)$ --- but its \emph{uniformity}. A decomposition
whose number of auxiliary (catalytic) variables grows with $n$ yields no
holonomicity statement. What is needed is a description of the constraint whose
exceptional data stays bounded as $n\to\infty$. That is what we supply.

\begin{theorem}\label{thm:main}
For every fixed $r,s\ge 1$, the sequences $\ars(n)$ and $\brs(n)$ are
holonomic.
\end{theorem}

Section~\ref{sec:bij} constructs the bijection and proves the seam count,
Section~\ref{sec:rep} records the resulting inclusion--exclusion
representation, and Section~\ref{sec:hol} assembles the closure argument.
Section~\ref{sec:ver} reports the computational verification.

\section{The interleaving bijection and its seams}\label{sec:bij}

\begin{definition}[Interleaving]
For $n\ge 1$ and $k\ge 1$ let $\iota_{n,k}$ be the sequence obtained by
listing $1,\dots,n$ in the order
\[
  1,\ 1+k,\ 1+2k,\ \dots,\quad 2,\ 2+k,\ 2+2k,\ \dots,\quad \dots,\quad
  k,\ 2k,\ 3k,\ \dots,
\]
that is, by concatenating the $k$ arithmetic progressions of step $k$ in
increasing order of their least element. We write $\iota_{n,k}(j)$ for the
$j$-th entry, so $\iota_{n,k}$ is a bijection $\{1,\dots,n\}\to\{1,\dots,n\}$.
\end{definition}

Interleaving turns a step of $k$ into a step of $1$, except where one
progression is exhausted and the next begins. We call those places
\emph{seams}.

\begin{lemma}[Seam count]\label{lem:seams}
Let $k\ge 1$ and $n\ge k$. The number of indices $j\in\{1,\dots,n-1\}$ with
$\iota_{n,k}(j+1)-\iota_{n,k}(j)\neq k$ is exactly $k-1$.
\end{lemma}

\begin{proof}
The list $\iota_{n,k}$ is the concatenation of the $k$ progressions
$P_c=\{c,c+k,c+2k,\dots\}\cap[1,n]$ for $c=1,\dots,k$, of lengths
\begin{equation}\label{eq:blocklen}
  |P_c| \;=\; \Bigl\lfloor \frac{n-c}{k} \Bigr\rfloor + 1 ,
  \qquad c=1,\dots,k .
\end{equation}
Within a single $P_c$ consecutive entries differ by exactly $k$, so no interior
index of a block is a seam. A seam can therefore occur only at an index $j$
where $P_c$ ends and $P_{c+1}$ begins, and there are exactly $k-1$ such
indices, one for each $c=1,\dots,k-1$. At such an index the difference is
$(c+1)-\max P_c$, which is negative and hence certainly not equal to $k$. Thus
every block boundary is a seam and there are $k-1$ of them, located at the
partial sums $|P_1|+\dots+|P_c|$, $c=1,\dots,k-1$.
\end{proof}

\begin{corollary}\label{cor:affine}
Fix $k$ and $c$. On each residue class of $n$ modulo $k$, the seam index
$|P_1|+\dots+|P_c|$ is an affine function of $n$ with slope $c/k$. Consequently,
for fixed $r,s$, all position and value seams are simultaneously affine in $n$
on each residue class of $n$ modulo $\operatorname{lcm}(r,s)$.
\end{corollary}

\begin{proof}
On a fixed residue class of $n$ mod $k$ the floor in \eqref{eq:blocklen} is
$(n-c-\rho_c)/k$ for a constant $\rho_c$, so each $|P_c|$ is affine in $n$ with
slope $1/k$, and a partial sum of $c$ of them is affine with slope $c/k$. The
position seams involve $k=r$ and the value seams $k=s$; both are affine on each
class modulo $\operatorname{lcm}(r,s)$.
\end{proof}

Now fix $r,s\ge 1$ and $n$. Define the position relabelling
$p=\iota_{n,r}$ and the value relabelling $v=\iota_{n,s}$, and associate to
$\pi\in\Sn$ the permutation $\sigma\in\Sn$ determined by
\begin{equation}\label{eq:transport}
  \sigma_k \;=\; v^{-1}\bigl(\pi_{\,p(k)}\bigr), \qquad k=1,\dots,n .
\end{equation}
Since $p$ and $v$ are bijections, \eqref{eq:transport} is a bijection
$\Sn\to\Sn$.

Write
\[
  \mathcal{P} = \{k : p(k+1)-p(k)\neq r\},\qquad
  \mathcal{V}^{+} = \{u : v(u+1)-v(u)\neq s\},\qquad
  \mathcal{V}^{-} = \{u : v(u-1)-v(u)\neq -s\}
\]
for the position seams and the forward and backward value seams. By
Lemma~\ref{lem:seams}, $|\mathcal{P}|=r-1$ and
$|\mathcal{V}^{+}|=|\mathcal{V}^{-}|=s-1$.

\begin{proposition}\label{prop:transport}
Under \eqref{eq:transport}:
\begin{enumerate}
\item[(i)] $\pi$ satisfies $\pi_{i+r}-\pi_i\neq s$ for all $i$ if and only if
$\sigma$ satisfies $\sigma_{k+1}-\sigma_k\neq 1$ for every $k$ with
$k\notin\mathcal{P}$ and $\sigma_k\notin\mathcal{V}^{+}$;
\item[(ii)] $\pi$ satisfies $|\pi_{i+r}-\pi_i|\neq s$ for all $i$ if and only if
$\sigma$ satisfies $|\sigma_{k+1}-\sigma_k|\neq 1$ for every $k\notin\mathcal{P}$,
where an ascent $\sigma_{k+1}=\sigma_k+1$ is excused when
$\sigma_k\in\mathcal{V}^{+}$ and a descent $\sigma_{k+1}=\sigma_k-1$ is excused
when $\sigma_k\in\mathcal{V}^{-}$.
\end{enumerate}
Consequently the exceptional data has size $(r-1)+(s-1)$ in case (i) and
$(r-1)+2(s-1)$ in case (ii); in particular it is bounded independently of $n$.
\end{proposition}

\begin{proof}
Consider a pair of adjacent new indices $k,k+1$ with $k\notin\mathcal{P}$. Then
$p(k+1)=p(k)+r$, so the pair corresponds to the old positions $i=p(k)$ and
$i+r$. Suppose further $\sigma_k=u\notin\mathcal{V}^{+}$, so $v(u+1)=v(u)+s$.
Then $\sigma_{k+1}=u+1$ if and only if
$\pi_{i+r}=v(u+1)=v(u)+s=\pi_i+s$, i.e.\ if and only if
$\pi_{i+r}-\pi_i=s$. This proves (i): outside the excused set the two
conditions agree pointwise, and on the excused set the adjacency condition is
imposed on neither side.

For (ii) the same computation applies to descents with $\mathcal{V}^{-}$ in
place of $\mathcal{V}^{+}$: with $k\notin\mathcal{P}$ and
$u\notin\mathcal{V}^{-}$ we have $v(u-1)=v(u)-s$, so $\sigma_{k+1}=u-1$ if and
only if $\pi_{i+r}-\pi_i=-s$. Since $|\pi_{i+r}-\pi_i|=s$ means
$\pi_{i+r}-\pi_i=\pm s$, both signs must be excused, and the exceptional data
consists of $\mathcal{P}$ together with $\mathcal{V}^{+}$ and $\mathcal{V}^{-}$.
The cardinalities follow from Lemma~\ref{lem:seams}.
\end{proof}

\begin{remark}
For $(r,s)=(2,2)$ the excused set has size $1+1=2$ and
Proposition~\ref{prop:transport}(i) recovers Matsuo's bijection as used in
\cite{SZ}, where the condition is stated as ``$\pi_{i+1}-\pi_i\neq 1$ except
when $i=\lfloor (n+1)/2\rfloor$ or $\pi_i=\lfloor (n+1)/2\rfloor$''.
\end{remark}

\section{The two-layer inclusion--exclusion representation}\label{sec:rep}

Proposition~\ref{prop:transport} reduces both families to a single counting
problem. Fix $n$, a set $\mathcal{P}\subseteq\{1,\dots,n-1\}$ of excused
positions and a set $\mathcal{V}\subseteq\{1,\dots,n\}$ of excused values, and
let
\[
  R(n;\mathcal{P},\mathcal{V})
  = \#\bigl\{\sigma\in\Sn :\ \sigma_{k+1}-\sigma_k\neq 1
    \ \text{ whenever } k\notin\mathcal{P} \text{ and } \sigma_k\notin\mathcal{V}\bigr\},
\]
with the analogous two-sign version $R^{\pm}$ in the $\brs$ case. We must show
that $R$ admits a representation whose number of catalytic variables is
governed by $|\mathcal{P}|$ and $|\mathcal{V}|$ alone.

\subsection*{Bonds and blocks}

Call a pair $(k,v)$ a \emph{violation} of $\sigma$ if $\sigma_k=v$,
$\sigma_{k+1}=v+1$, $k\notin\mathcal{P}$ and $v\notin\mathcal{V}$. Thus
$R(n;\mathcal{P},\mathcal{V})$ counts permutations with no violation, and
inclusion--exclusion over sets of violations gives
\begin{equation}\label{eq:ie1}
  R(n;\mathcal{P},\mathcal{V})
  \;=\; \sum_{S} (-1)^{|S|}\, A(S),
\end{equation}
where $S$ ranges over sets of \emph{value bonds}
$S\subseteq\{1,\dots,n-1\}\setminus\mathcal{V}$ and $A(S)$ counts the
permutations in which every $v\in S$ occurs as a violation.

Forcing the bonds in $S$ glues the values into maximal runs: if $S$ decomposes
$\{1,\dots,n\}$ into $n-|S|$ maximal runs of consecutive integers, then a
permutation realising all bonds of $S$ is exactly an arrangement of these
$n-|S|$ \emph{blocks} in some order, read left to right. Without the position
constraint one would obtain the classical identity
$\sum_S(-1)^{|S|}(n-|S|)!$; the position constraint $k\notin\mathcal{P}$ is
what makes the problem two-layered.

\subsection*{The position layer}

In an arrangement of blocks, the internal bonds occupy the positions
determined by the partial sums of the block sizes: if the blocks have sizes
$\ell_1,\dots,\ell_{n-|S|}$ in the order used, the internal bonds of the $j$-th
block occupy positions
\[
  L_{j-1}+1,\ \dots,\ L_{j}-1,
  \qquad L_j=\ell_1+\dots+\ell_j .
\]
$A(S)$ therefore counts those arrangements in which no internal bond position
lies in $\mathcal{P}$. Since $|\mathcal{P}|=r-1$ is bounded, we may remove this
constraint by a second, \emph{finite}, inclusion--exclusion:
\begin{equation}\label{eq:ie2}
  A(S) \;=\; \sum_{T\subseteq\mathcal{P}} (-1)^{|T|}\, B(S,T),
\end{equation}
where $B(S,T)$ counts arrangements of the blocks of $S$ in which every
position of $T$ \emph{is} an internal bond position (and the positions of
$\mathcal{P}\setminus T$ are unconstrained). The sum \eqref{eq:ie2} has
$2^{|\mathcal{P}|}=2^{\,r-1}$ terms, a number depending only on $r$.

\subsection*{Block factors}

A block of size $\ell$ absorbs $\ell-1$ bonds and occupies $\ell$ consecutive
positions, so with $x$ marking positions the unmarked block factor is
\begin{equation}\label{eq:blockfactor}
  \mathcal{B}(x)\;=\;\sum_{\ell\ge 1}(-1)^{\ell-1}x^{\ell}\;=\;\frac{x}{1+x} .
\end{equation}
In the two-sign case a block of size $\ell\ge 2$ may be traversed in either
direction --- the run $v,v+1,\dots$ may appear increasing or decreasing --- and
a block of size $1$ in only one way, giving
\begin{equation}\label{eq:blockfactorpm}
  \mathcal{B}^{\pm}(x)\;=\;x+2\sum_{\ell\ge 2}(-1)^{\ell-1}x^{\ell}
  \;=\;\frac{x-x^{2}}{1+x},
\end{equation}
which is the block generating function of the Hertzsprung problem.

\begin{remark}
As a check on \eqref{eq:blockfactor} and on the umbral normalisation below,
take $\mathcal{P}=\mathcal{V}=\emptyset$ and $n=3$. Then
$[x^{3}]\mathcal{B}^{j}=1,-2,1$ for $j=1,2,3$, and
$1\cdot 1!-2\cdot 2!+1\cdot 3!=3$, which is the number of permutations of
$\{1,2,3\}$ with no succession, namely $132$, $213$, $321$.
\end{remark}

\subsection*{Where the two kinds of mark act}

The two excusal types enter in genuinely different ways, and it is worth
separating them.

\emph{Excused values force block boundaries.} If $v\in\mathcal{V}$ then the
bond $v$ is never a violation, so $v\notin S$ for every $S$ occurring in
\eqref{eq:ie1}. Hence no block may span the gap between $v$ and $v+1$: the
excused values cut $\{1,\dots,n\}$ into $|\mathcal{V}|+1$ \emph{value
segments}, and the generating function factorises as a product of one block
sequence per segment. A catalytic variable $y_i$ marking the $i$-th cut records
its location.

\emph{Excused positions split the arrangement.} The position of a bond is
determined by the \emph{arrangement} of the blocks, not by their value order,
so a position mark cannot be carried inside a block factor. Instead, fixing
which arrangement positions $t_1<\dots<t_q$ are internal bond positions cuts
the arrangement into $q+1$ consecutive \emph{arrangement segments}. Blocks may
be distributed among the segments freely and ordered arbitrarily within each,
so if the segments receive $j_0,\dots,j_q$ blocks the number of orderings is
$j_0!\,j_1!\cdots j_q!$. This is realised by giving each segment its own
block-counting variable $z_0,\dots,z_q$ and applying the multivariate umbral
functional
\begin{equation}\label{eq:umbral}
  \bigl\langle\,\cdot\,\bigr\rangle:\quad
  z_0^{\,j_0}\cdots z_q^{\,j_q}\;\longmapsto\; j_0!\cdots j_q!,
  \qquad
  \langle f\rangle=\int_0^\infty\!\!\cdots\!\int_0^\infty
  f\,e^{-z_0-\dots-z_q}\,dz_0\cdots dz_q .
\end{equation}

\begin{lemma}\label{lem:marked}
Let $\mathcal{B}^{\ast}(x;u)=\sum_{\ell\ge 2}(-1)^{\ell-1}x^{\ell}
\bigl(u+u^{2}+\dots+u^{\ell-1}\bigr)$ be the generating function of a block
carrying one distinguished interior slot, the variable $u$ recording the slot's
offset from the start of the block. Then
\[
  \mathcal{B}^{\ast}(x;u)\;=\;\frac{-x^{2}u}{(1+x)(1+xu)} .
\]
\end{lemma}

\begin{proof}
Summing the inner geometric series,
$u+\dots+u^{\ell-1}=u(u^{\ell-1}-1)/(u-1)$, so
\[
  \mathcal{B}^{\ast}
  =\frac{u}{u-1}\Bigl(\tfrac1u\sum_{\ell\ge2}(-1)^{\ell-1}(xu)^{\ell}
   -\sum_{\ell\ge2}(-1)^{\ell-1}x^{\ell}\Bigr).
\]
For any $y$, $\sum_{\ell\ge1}(-1)^{\ell-1}y^{\ell}=y/(1+y)$, whence
$\sum_{\ell\ge2}(-1)^{\ell-1}y^{\ell}=y/(1+y)-y=-y^{2}/(1+y)$. Substituting
$y=xu$ and $y=x$,
\[
  \mathcal{B}^{\ast}
  =\frac{u}{u-1}\Bigl(\frac{-x^{2}u}{1+xu}+\frac{x^{2}}{1+x}\Bigr)
  =\frac{ux^{2}}{u-1}\cdot\frac{(1+xu)-u(1+x)}{(1+x)(1+xu)}
  =\frac{ux^{2}}{u-1}\cdot\frac{1-u}{(1+x)(1+xu)},
\]
since $(1+xu)-u(1+x)=1-u$. The factor $(1-u)/(u-1)=-1$ gives the claim.
\end{proof}

\begin{remark}
For $(r,s)=(2,2)$ there is one position seam and one value seam, so
\eqref{eq:umbral} has two variables and the value line is cut into two
segments. This is exactly the structure of the representation used in
\cite{PadhiA22}, where the umbral functional is $z^{k}w^{j}\mapsto k!\,j!$ and
the marked-block factor is the $\mathcal{B}^{\ast}$ of
Lemma~\ref{lem:marked}.
\end{remark}

Assembling \eqref{eq:ie1}, \eqref{eq:ie2}, \eqref{eq:blockfactor} and
\eqref{eq:umbral}:

\begin{proposition}\label{prop:rep}
For every $n$, $\mathcal{P}$ and $\mathcal{V}$,
\begin{equation}\label{eq:rep}
  R(n;\mathcal{P},\mathcal{V})
  \;=\;
  \sum_{T\subseteq\mathcal{P}}(-1)^{|T|}\,
  \bigl[x^{n}\,y_1^{v_1}\cdots y_p^{v_p}\bigr]\,
  \bigl\langle\,\Phi_T\bigl(x,z_0,\dots,z_{|T|},y_1,\dots,y_p\bigr)\bigr\rangle ,
\end{equation}
where $p=|\mathcal{V}|$, the functional $\langle\cdot\rangle$ is
\eqref{eq:umbral}, and each $\Phi_T$ is a rational function built from
\eqref{eq:blockfactor} --- respectively \eqref{eq:blockfactorpm} in the
two-sign case --- as a product of $p+1$ value-segment factors distributed over
$|T|+1$ arrangement segments. In particular $\Phi_T$ involves
$|T|+1\le r$ umbral variables and $p$ catalytic variables.
\end{proposition}

\begin{corollary}\label{cor:bounded}
For fixed $r,s$ the representation \eqref{eq:rep} is a sum of at most
$2^{\,r-1}$ terms, each a rational function in at most $1+r$ umbral variables
together with at most $s-1$ catalytic variables in the $\ars$ case and
$2(s-1)$ in the $\brs$ case. All of these counts are independent of $n$.
\end{corollary}

\begin{proof}
Immediate from Proposition~\ref{prop:rep} and the seam counts of
Proposition~\ref{prop:transport}, since $|\mathcal{P}|=r-1$ and
$|\mathcal{V}|=s-1$ (resp.\ $2(s-1)$).
\end{proof}

\section{Holonomicity}\label{sec:hol}

\begin{proof}[Proof of Theorem~\ref{thm:main}]
Fix $r,s$. By Proposition~\ref{prop:transport} and
Corollary~\ref{cor:bounded}, $\ars(n)$ (resp.\ $\brs(n)$) is obtained from a
rational function $F$ in a fixed finite set of variables --- one group of tile
variables, and $c\le (r-1)+2(s-1)$ catalytic variables --- by the following
finite sequence of operations, each of which preserves holonomicity:
\begin{enumerate}
\item[(1)] \emph{Coefficient extraction.} Extracting the coefficient of
$x^{n}y_1^{v_1}\cdots y_p^{v_p}$ from a rational function in finitely many
variables, and more generally taking a diagonal, yields a D-finite series
(Lipshitz \cite{Lipshitz}; Christol). By Corollary~\ref{cor:bounded} the number
of extractions is at most $1+p$, independent of $n$.
\item[(2)] \emph{The umbral functional.} Applying $z^{j}\mapsto j!$ in a
variable is integration against $e^{-z}$, and the resulting parameterised
integral of a D-finite integrand is D-finite in the remaining parameters
(Almkvist--Zeilberger \cite{AZ}; Lipshitz \cite{Lipshitz} for the closure
statement). By Corollary~\ref{cor:bounded} the functional \eqref{eq:umbral}
involves at most $1+r$ variables, again independent of $n$.
\item[(3)] \emph{Specialisation of the excused data.} By
Corollary~\ref{cor:affine} the seam indices are affine in $n$ on each residue
class of $n$ modulo $L:=\operatorname{lcm}(r,s)$, with slopes $c/r$ and $c/s$.
These slopes are not integers, so we reparametrise: fix a residue
$\rho\in\{0,\dots,L-1\}$ and write $n=LM+\rho$. Since $r\mid L$ and $s\mid L$,
each seam index becomes an affine function of $M$ with \emph{integer} slope and
integer constant term. Substituting integer-affine functions of $M$ for the
catalytic exponents of a holonomic multivariate sequence is a diagonal-type
specialisation, and diagonals of D-finite series are D-finite
(Lipshitz \cite{LipshitzDiag}); see also \cite[\S 2]{Zeilberger} for the
holonomic-systems formulation. Hence each of the $L$ sections
$M\mapsto \ars(LM+\rho)$ is holonomic.
\item[(4)] \emph{Interlacing.} Finally $\ars$ (resp.\ $\brs$) is the
interlacing of the $L$ sections produced in (3), and an interlacing of finitely
many holonomic sequences is holonomic (Kauers--Paule
\cite[Ch.~7]{KauersPaule}).
\end{enumerate}
Each step preserves holonomicity and the number of steps depends only on
$(r,s)$, not on $n$. Hence $\ars$ and $\brs$ are holonomic.
\end{proof}

\begin{remark}
The argument gives no bound on the order or degree of the resulting recurrence,
and is not effective in any practical sense: for $(r,s)=(2,2)$ the minimal
operator has order $8$ and degree $11$ \cite{KK}, and obtaining it explicitly is
a substantial computation independent of the above. Challenge 3 asks only for
holonomicity, which is what is proved here.
\end{remark}

\section{Computational verification}\label{sec:ver}

Proposition~\ref{prop:transport} was verified by brute force for all
$4\le n\le 8$ and the pairs listed below: for each, the permutations of $\Sn$
were enumerated, $\ars(n)$ and $\brs(n)$ computed directly from the
definitions, the relabellings $p,v$ constructed, and the transformed counts
compared. Every count agreed exactly, and the seam cardinalities were as
predicted by Lemma~\ref{lem:seams} in every case.

\begin{center}
\begin{tabular}{lccc}
\toprule
$(r,s)$ & excused, $\ars$ & excused, $\brs$ & agreement, $4\le n\le 8$\\
\midrule
$(2,2)$ & $2$ & $3$ & exact \\
$(2,3)$ & $3$ & $5$ & exact \\
$(2,4)$ & $4$ & $7$ & exact \\
$(3,2)$ & $3$ & $4$ & exact \\
$(3,3)$ & $4$ & $6$ & exact \\
$(3,4)$ & $5$ & $8$ & exact \\
$(4,2)$ & $4$ & $5$ & exact \\
$(4,3)$ & $5$ & $7$ & exact \\
$(4,4)$ & $6$ & $9$ & exact \\
\bottomrule
\end{tabular}
\end{center}

For instance $a_{2,2}(4,\dots,8) = 18, 75, 410, 2729, 20906$ and
$b_{2,2}(4,\dots,8) = 16, 44, 200, 1288, 9512$, the latter matching
\texttt{A110128}. In the $\brs$ case the naive one-sided excusal rule
undercounts and a strictly larger rule overcounts; only the symmetrised rule of
Proposition~\ref{prop:transport}(ii) reproduces $\brs$, which is the
computational signature of both signs needing to be excused.

\section*{Acknowledgements}

The problem is Challenge 3 of Spahn and Zeilberger \cite{SZ}; the $(2,2)$
bijection on which the construction is modelled is due to Matsuo.

\begin{thebibliography}{9}

\bibitem{AZ} G.~Almkvist and D.~Zeilberger, \emph{The method of differentiating
under the integral sign}, J. Symbolic Comput. \textbf{10} (1990), 571--591.

\bibitem{KauersPaule} M.~Kauers and P.~Paule, \emph{The Concrete Tetrahedron},
Springer, 2011.

\bibitem{KK} M.~Kauers and C.~Koutschan, \emph{Guessing with little data},
arXiv:2202.07966.

\bibitem{LipshitzDiag} L.~Lipshitz, \emph{The diagonal of a D-finite power
series is D-finite}, J. Algebra \textbf{113} (1988), 373--378.

\bibitem{Lipshitz} L.~Lipshitz, \emph{D-finite power series}, J. Algebra
\textbf{122} (1989), 353--373.

\bibitem{PadhiA22} J.~S.~Padhi, \emph{An integral representation for
$a_{2,2}$}, manuscript in preparation. (Cited only for the explicit
$(r,s)=(2,2)$ specialisation in the remarks; no result of the present paper
depends on it.)

\bibitem{SZ} E.~Spahn and D.~Zeilberger, \emph{Two-dimensional restricted
permutations}, arXiv:2211.02550; also ECA \textbf{3} (2023), S2A10.

\bibitem{Zeilberger} D.~Zeilberger, \emph{A holonomic systems approach to
special functions identities}, J. Comput. Appl. Math. \textbf{32} (1990),
321--368.

\end{thebibliography}

\end{document}
